\documentclass[11pt]{article}
\usepackage{mathblatt}


\blattfuss{Gleichungssysteme}{Lösungen}
\begin{document}
\noindent{\large\bfseries Gleichungssysteme \textperiodcentered{} Lösungen}\par\medskip
\begin{pfloesung}{}
\lz{1.}{$15 - x$}{$y$}{}
\end{pfloesung}
\begin{pfloesung}{}
\lz{2.}{Gerade durch $(0|1)$ und $(1|3)$}{}{}
\end{pfloesung}
\begin{pfloesung}{}
\lz{3.}{$6$}{$x$}{}
\end{pfloesung}
\begin{pfloesung}{}
\lz{4.}{$3x + 12$}{}{}
\end{pfloesung}
\begin{pfloesung}{}
\lz{5.}{$7{,}20\,€$}{}{}
\end{pfloesung}
\begin{pfloesung}{}
\lz{6.}{$1{,}50\,€$}{$3x = 4{,}50$; $x$}{}
\end{pfloesung}
\begin{pfloesung}{}
\lz{7.}{von beiden}{}{}
\end{pfloesung}
\begin{pfloesung}{}
\lz{8.}{vor $x$ und $y$: $2$, $5$ und $3$, $1$}{rechts: $58\,€$ und $35\,€$; in €}{}
\end{pfloesung}
\begin{pfloesung}{}
\lz{9.}{$y$: $5$, $4$, $3$, $2$}{}{}
\end{pfloesung}
\begin{pfloesung}{}
\lz{10.}{$69$}{I: $x + y = 21$, II: $2x + 5y$}{}
\end{pfloesung}
\begin{pfloesung}{}
\lz{11.}{$68$}{I: $x + y = 23$, II: $2x + 4y$}{}
\end{pfloesung}
\begin{pfloesung}{}
\lz{12.}{$7$}{I: $x + y = 31$, II: $x - y$}{}
\end{pfloesung}
\begin{pfloesung}{}
\lz{13.}{$28$}{I: $x + y = 42$, II: $x = 2y$; $y = 14$, $x$}{}
\end{pfloesung}
\begin{pfloesung}{{\small\color{mbgrau}P10 2016 · 6d}}
\lz{14.}{7 Drei-Bett-Zimmer; 9 Fünf-Bett-Zimmer}{x + y = 16; 3x + 5y = 66; 3(16 $-$ y) + 5y = 66 $\Rightarrow$ y = 9}{}
\end{pfloesung}
\begin{pfloesung}{}
\lz{15.}{$118{,}40$}{I: $x + y = 120$, II: $0{,}80x + 1{,}20y$}{}
\end{pfloesung}
\begin{pfloesung}{{\small\color{mbgrau}P10 2021 · 7b}}
\lz{16.}{I: 2x + 2y = 77,80, II: x + 3y = 64,90; Erwachsener 25,90 €, Kind 13,00 €}{x = 64,90 $-$ 3y; 129,80 $-$ 4y = 77,80 $\Rightarrow$ y = 13}{}
\end{pfloesung}
\begin{pfloesung}{}
\lz{17.}{Gegenzahlen}{}{}
\end{pfloesung}
\begin{pfloesung}{}
\lz{18.}{Gleichung II}{}{}
\end{pfloesung}
\begin{pfloesung}{}
\lz{19.}{Lösung $(3|7)$}{$x = 3$, $y = 7$}{}
\end{pfloesung}
\begin{pfloesung}{}
\lz{20.}{Lösung $(3|5)$}{$x = 3$, $y = 5$}{}
\end{pfloesung}
\begin{pfloesung}{}
\lz{21.}{Lösung $(4|1)$}{$x = 4$, $y = 1$}{}
\end{pfloesung}
\begin{pfloesung}{}
\lz{22.}{Lösung $(1|3)$}{$y = 3$, $x = 1$}{}
\end{pfloesung}
\begin{pfloesung}{}
\lz{23.}{Lösung $(3|5)$}{$x = 3$, $y = 5$}{}
\end{pfloesung}
\begin{pfloesung}{}
\lz{24.}{Lösung $(5|9)$}{$x = 5$, $y = 9$}{}
\end{pfloesung}
\begin{pfloesung}{}
\lz{25.}{Lösung $(-2|3)$}{$x = -2$, $y = 3$}{}
\end{pfloesung}
\begin{pfloesung}{}
\lz{26.}{Lösung $(3|5)$}{$x = 3$, $y = 5$}{}
\end{pfloesung}
\begin{pfloesung}{}
\lz{27.}{Lösung $(5|2)$}{$x = 5$, $y = 2$}{}
\end{pfloesung}
\begin{pfloesung}{}
\lz{28.}{Lösung $(-1|4)$}{$x = -1$, $y = 4$}{}
\end{pfloesung}
\begin{pfloesung}{}
\lz{29.}{Lösung $(3|1)$}{$x = 3$, $y = 1$}{}
\end{pfloesung}
\begin{pfloesung}{}
\lz{30.}{Lösung $(5|2)$}{$x = 5$, $y = 2$}{}
\end{pfloesung}
\begin{pfloesung}{}
\lz{31.}{$20{,}60$}{$y = 54 - x$ in II: $7x + 216 - 4x = 316{,}20$, $x = 33{,}40$; $y$}{}
\end{pfloesung}
\begin{pfloesung}{}
\lz{32.}{Lösung $(1|5{,}5)$}{I: $y = 4x + 1{,}5$; in II: $5x + 8x + 3 = 16$, $x = 1$; $y = 5{,}5$}{}
\end{pfloesung}
\begin{pfloesung}{{\small\color{mbgrau}P10 2022 · 7b}}
\lz{33.}{Apfelbaum 27,20 €, Pflaumenbaum 35,80 €}{y = 63 $-$ x; 6x = 163,20}{}
\end{pfloesung}
\begin{pfloesung}{}
\lz{34.}{$x$: Preis eines Blocks in €, $y$: Preis eines Füllers in €}{}{}
\end{pfloesung}
\begin{pfloesung}{}
\lz{35.}{$x$: Preis eines Zelts in €, $y$: Preis eines Schlafsacks in €}{}{}
\end{pfloesung}
\begin{pfloesung}{}
\lz{36.}{$230$: Einnahmen zusammen in €}{$8$: Preis eines großen Glases in €}{}
\end{pfloesung}
\begin{pfloesung}{{\small\color{mbgrau}P10 2022 · 7a}}
\lz{37.}{x: Preis eines Apfelbaums in €, y: Preis eines Pflaumenbaums in €}{Koeffizienten 9 und 3 = Stückzahlen}{}
\end{pfloesung}
\begin{pfloesung}{}
\lz{38.}{Im Korb liegen zusammen $25$ Äpfel und Birnen}{}{}
\end{pfloesung}
\begin{pfloesung}{}
\lz{39.}{$x$: Preis einer Tischtennisplatte in €, $y$: Preis eines Kickers in €}{}{}
\end{pfloesung}
\begin{pfloesung}{\textbf{Prüfstein} \textperiodcentered{} P10 2024 · 7}
\lz{a)}{I. x + y = 3,80; II. 6x + 5y = 21,20}{x Preis Rose, y Preis Tulpe}{}
\lz{b)}{Filip kauft insgesamt 13 Blumen (Rosen und Tulpen zusammen); r = 8 Rosen, t = 5 Tulpen}{r = 13 $-$ t; 29,9 $-$ 0,6t = 26,9 $\Rightarrow$ 0,6t = 3}{}
\end{pfloesung}
\end{document}